\chapter{Gauss Norm}

Thus, far we have shown that for a non-archimedean normed ring $R$, $R_c \langle x_1, \dots, x_n \rangle$
is a ring for any parameter $c$. We now want to endow this with a non-archimedea absolute value.

Towards this let us consider the following function on multivariate power series:

\begin{definition}
Define $\lvert \cdot \rvert _c :  R_c \langle x_1, \dots, x_n \rangle \to \R$ where
\[
    \lvert f \rvert _c = \sup_{i \in \mathbb{N}^n} \lvert a_{(i_1, \cdots, i_n)} \rvert
    \prod_n \lvert c \rvert _n^{i_n},
\]
for $f(x) = \sum a_{(i_1, \cdots, i_n)} x_1^{i_1} \cdots x_n^{i_n}$ and call this the Gauss norm
with parameter $c$.
\end{definition}

We may think to ask is this a norm on the whole space of multivariate power series; but this
has a very easy answer. Simply, for a general power series
\[
  \sup_{i \in \mathbb{N}^n} \lvert a_{(i_1, \cdots, i_n)} \rvert
    \prod_n \lvert c \rvert _n^{i_n}
\]
need not exist - specifically, unless the set is bounded above this would be undefined in $\R$.

Moreover, even if we restrict to power series that have bounded
\[
  \lvert a_{(i_1, \cdots, i_n)} \rvert \prod_n \lvert c \rvert _n^{i_n},
\]
which we will refer to as "having gauss norm". A product of two such power series need not have gauss
norm.
For example, consider $c = 1$ and
\[
  f = \sum x^n \in \Q_p [ \! [x] \! ]
\]
then
\[
  f^2 = \sum (n + 1) x^n,
\]
and we have $\lvert n + 1 \rvert_p \tendsto \infty$.

Nevertheless, restricting to $R_c \langle x_1, \dots, x_n \rangle$ is sufficient to prove what we want.

\begin{theorem}
  For $n$ a natural number, $c$ a tuple of non-zero real numbers and $R$ a non-archimedean ring. Then
  the gauss norm is a non-archimedean absolute value on $R_c \langle x_1, \dots, x_n \rangle$.
\end{theorem}

\begin{remark}
  We could alternatively define the gauss norm without $\lvert \cdot \rvert$ around the $c$ and instead
  require that $c$ is strictly positive in the theorem. The latter is the choice we make in Lean; this is
  to match the definition of restricted power series - where we do not have norms around the $c$ either -
  and to prevent having $\lvert \cdot \rvert_a = \lvert \cdot \rvert_{-a}.$

  We could argue to also require $0 \leq c$ in the definition of restricted power series as well, but in
  interest of keeping definitions the most general we opt to allow negative entries in $c$.
\end{remark}

Note we really require no entries of $c$ to be zero, as it is immediate to see that in this case we
would have $\lvert \cdot \rvert_c \equiv 0$.

\section{Lean implementation}

To prove the theorem, for each property of an absolute value we actually generalise the conditions
on $R$ and $c$ as much as possible. For example, instead of $R$ being a normed ring we only impose the
following variables:

\begin{itemize}
  \item $R$ is a semiring; and
  \item $v$ is a function $R \to \R$.
\end{itemize}

Then when required we strengthen can strengthen $R$ to a ring and impose properties on $v$.

Moreover, although we established that the gauss norm is not a norm on the space of power series -
some properties can still be proved, and assuming the power series "has gauss norm" we can prove
everything except multiplicativity.

We will now state the harder statements we give towards showing the gauss norm is a
non-archimedean absolute value on $R_c \langle x_1, \dots, x_n \rangle$.

First, we generalise the definition of gauss norm into and formally define what we mean by a power
series "has gauss norm".

\begin{definition}
  We define the gauss norm of a power series $f \in R_c [ \! [ x_1, \dots, x_n ] \! ]$ as
  \[
    \lvert f \rvert_c = \sup_{(i_1, \cdots, i_n)} v (a_{(i_1, \cdots, i_n)})  \prod_n \lvert c \rvert _n^{i_n}
  \]
  where $f = \sum a_{(i_1, \cdots, i_n)} x_1^{i_1} \cdots x_n^{i_n}$.
\end{definition}

\begin{definition}
  Let $R$ be a semiring. We say a power series $f = \sum a_{(i_1, \cdots, i_n)} x_1^{i_1} \cdots x_n^{i_n}$
  has gauss norm if the set
  \[
    v (a_{(i_1, \cdots, i_n)})  \prod_n \lvert c \rvert _n^{i_n}
  \]
  is bounded above.
\end{definition}

Then, the non-archimedean property can be stated as.

\begin{lemma}
  If $R$ is a semiring, $c$ is non-negative and $v$ is a non-negative non-archimedean function then
  for two power series $f, g \in R_c [ \! [ x_1, \dots, x_n ] \! ]$ that have gauss norm, then
  \[
    \lvert f + g \rvert_c \leq \max \left\{ \lvert f \rvert_c, \lvert g \rvert_c \right\}.
  \]
\end{lemma}

%% I also want to state the gaussNorm_mul_le and gaussNorm_le_mul + antidiagonal_dominant statements
%% But need to decide how I want to do this
